\documentclass[a4paper,12pt]{report}
\newcommand{\path}{../}
\usepackage[fiche,niveau={1MA1},nfiche={18}, annee={Emilie-Gourd, 2024--2025},
auteur={ns},theme={Fonctions -- Géométrie analytique (introduction)}]{\path packages/bandeaux}
\usepackage{\path packages/boites}
\usepackage{\path packages/fontes}
\usepackage{\path packages/courslatex}
\usepackage{\path packages/mathenv}
\allowdisplaybreaks

\begin{document}
{\bfseries Compétences :}
\begin{itemize}
    \item Calculer la distance entre deux points
    \item Déterminer le milieu d'un segment   
    \item Résoudre des problèmes à l'aide des notions étudiées
\end{itemize}
\begin{defi}[distance entre deux points]
La {\bfseries distance entre deux points} $A(x_1;y_1)$ et $B(x_2;y_2)$ est donnée par la formule \[\text{dist}(A;B)=\sqrt{(x_1-x_2)^2 +(y_1-y_2)^2 }.\]
\end{defi}
    
\begin{rem}
La formule de la distance entre deux points est une application directe du théorème de Pythagore.
\begin{center}
    \begin{tikzpicture}[scale=0.7]
\begin{axis}[
    label style={font=\scriptsize},
    tick label style={font=\scriptsize},
    axis lines=middle,
    scale only axis,                % Ensure only the axis is scaled
    axis equal image=true,          % Maintain equal scaling on both axes
    grid=major,
    xmin=0, xmax=7,
    ymin=0, ymax=5,
    xtick={0,1,2,3,4,5,6,7},
    ytick={0,1,2,3,4,5},
    xlabel={$x$},
    ylabel={$y$}
]

% Define points A, B, and C using axis coordinates
\coordinate (A) at (axis cs:2,1);
\coordinate (B) at (axis cs:6,4);
\coordinate (C) at (axis cs:2,4);

% Plot points A and B
\addplot[only marks, mark=*] coordinates {(2,1) (6,4)};

% Draw the right triangle sides using axis coordinates
\draw[blue] (A) -- (C) -- (B);

% Draw the hypotenuse (distance line) as a dashed line
\draw[blue, dashed] (A) -- (B);

% Labels for points A, B, and C
\node[below left] at (A) {\scriptsize$A(2,1)$};
\node[above right] at (B) {\scriptsize$B(6,4)$};
\node[left]       at (C) {\scriptsize$C(2,4)$};

% Horizontal distance label between A and C
\draw[<-] (axis cs:2,1.2) -- (axis cs:2,3.8) node[midway, left] {\scriptsize$3$};

% Vertical distance label between C and B
\draw[<-] (axis cs:3,4) -- (axis cs:5,4) node[midway, above] {\scriptsize$4$};

% Label the distance formula near the hypotenuse
\node[below right, sloped] at ($(A)!0.2!(B)$) 
    {\scriptsize$d = \sqrt{(6-2)^2 + (4-1)^2} = 5$};

\end{axis}
\end{tikzpicture}
\end{center}
\end{rem}

\begin{boiteExT}[Distance entre $A(2;-3)$ et $B(-4;5)$]
    \vspace{4cm}
\end{boiteExT}

\begin{defi}[milieu d'un segment]
    Le {\bfseries milieu d'un segment} $[AB]$ est le point $M$ tel que
    \[\text{dist}(A;M)=\text{dist}(M;B).\]
\end{defi}
        
\begin{rem}
    Si $A(x_1;y_1)$ et $B(x_2;y_2)$ sont deux points, le milieu $M$ du segment
    $[AB]$ est donné par les coordonnées
\[M=\left(\dfrac{x_1+x_2}{2};\dfrac{y_1+y_2}{2}\right).\]
\end{rem}
    
\begin{boiteExT}[Milieu du segment $[AB]$ avec $A(2;-3)$ et $B(-4;5)$]
    \vspace{4cm}
\end{boiteExT}

\begin{exo}
\begin{tasks}
    \task Déterminer la distance entre les points $A$ et $B$, et entre les points $C$ et $D$ avec $A(1;3)$, $B(4;9)$, $C(-3;1)$ et $D(-2;-4)$. 
    \task Déterminer le coordonées du milieu du segment $[AB]$ et du segment $[CD]$ avec $A, B, C, D$ donnés ci-dessus. 
\end{tasks}
\end{exo}


    
\begin{defi}[hauteur issue d'un sommet]
    Soit le triangle $ABC$. La {\bfseries hauteur issue du sommet} $A$ est la droite passant par $A$ perpendiculaire au côté opposé à $A$ (ici au côté $BC$). Par abus de langage, la longueur de la hauteur est la longueur du segment $[AH]$ où $H$ est la point d'intersection entre de la hauteur et le côté oppoés à $A$. 
\end{defi}

\begin{rem}
Soit le triangle $ABC$. Pour déterminer la hauteur issue du sommet $A$, on détermine l'équation de la droite du côté $BC$. Puis on détermine l'équation de la droite passant par $A$ et perpendiculaire à la droite $BC$.

Pour déterminer la longueur de la hauteur, on détermine le point d'intersection entre la hauteur et $BC$, puis on détermine la distance séparant $A$ et ce point d'intersection. 
\end{rem}
   \begin{boiteExT}[Déterminer la hauteur issue de $A$]
    \vspace{14cm}
       
   \end{boiteExT}
\medskip
   \begin{exo}
       \begin{tasks}
        \task Déterminer la longueur de la hauteur du triangle $ABC$ issue de $A$ où $A(0;1)$, $B(-2;3)$ et $C(-8;7)$.
        \task Choisir trois points non colinéaires dans le plan, calculer l'aire du triangle définit par ces points. 
       \end{tasks}
   \end{exo}
\end{document}

        \task Prouver que le quadrilatère $ABCD$ avec $A(0;2), B(1;0), C(3;1)$ et $D(2;3)$ est un carré.
En utilisant les propriétés des quadrilatères, on peut déterminer si un quadrilatère défini par quatre points est un carré ou un rectangle.

\begin{rem}
Il y a plusieurs possibilités pour déterminer si quatre points forment un
carré, en voici une. On vérifie que les quatre côtés ont le même longueur et qu'un des côté est perpendiculaire à un autre (ce qui permet de s'assurer que les quatre autres angles sont des angles droits). 
    
        En utilisant un raisonnement similaire, on peut déterminer si un quadrilatère en un rectangle. 
\end{rem}

\begin{boiteExT}[Déterminer si un quadrilatère est un carré]
    \vspace{15cm} 
\end{boiteExT}


